% !TEX root = ../main.tex
\chapter{PRISM: o visualizador de RTL}
\label{cap:prism}

O \tool{PRISM} (\emph{Processor Rendering Interface for Schematic Models}) desenha o seu projeto como um \textbf{diagrama de circuito navegável}: módulos, portas e conexões, do top-level até dentro do processador. É a resposta visual à pergunta \enquote{que hardware o meu código virou?}.

\figurapendente{Janela do PRISM com o diagrama do projeto}{O top-level do \code{MeuFiltro} renderizado; se possível, um segundo print já dentro do módulo do processador.}

\section{Como abrir}

Botão \botao{PRISM} na toolbar (requer um \textbf{top-level} definido no projeto --- Seção~\ref{sec:arquivos-do-projeto}). A AURORA então:

\begin{enumerate}
  \item compila o que estiver pendente e verifica a sintaxe dos Verilog;
  \item reúne todos os fontes do projeto (HDL do SAPHO, sintetizáveis, top-level, \file{Hardware/} de cada processador);
  \item sintetiza com o \tool{Yosys} em modo diagrama (sem mapeamento tecnológico --- o desenho preserva a estrutura RTL do seu código);
  \item renderiza cada módulo em SVG (via \tool{netlistsvg}) e abre a janela do PRISM.
\end{enumerate}

A saída do Yosys corre no terminal TVERI. A síntese também alimenta a \textbf{visão Hierarquia} da árvore de arquivos (Seção~\ref{sec:tres-arvores}).

\section{Navegando no diagrama}

A janela do PRISM é independente da principal (1400$\times$900, maximizada ao abrir), com o mesmo tema escuro:

\begin{itemize}
  \item \textbf{pan e zoom} com o mouse;
  \item \textbf{clique} num módulo entra nele --- o diagrama do submódulo já está materializado, a navegação é instantânea;
  \item \emph{breadcrumbs} no topo mostram o caminho na hierarquia e permitem voltar;
  \item \textbf{duplo clique} numa célula abre o \textbf{código-fonte} correspondente no editor principal, na linha exata da declaração --- do desenho ao código num gesto;
  \item \botao{Recompile} refaz a síntese e o diagrama após mudanças.
\end{itemize}

Os blocos do processador SAPHO usam \emph{skins} próprias --- símbolos desenhados para a ULA, memórias, pilhas --- em vez de retângulos genéricos.

\section{Modo simulação (DigitalJS)}

Além do diagrama estático, o PRISM tem um modo \textbf{Simular}: o circuito é convertido para o \tool{DigitalJS} e vira um esquemático \emph{interativo}, onde você alterna entradas e vê os valores se propagarem pelas conexões, ao vivo. É uma lupa didática para circuitos pequenos:

\begin{itemize}
  \item o limite é de \textbf{3000 células} --- acima disso o PRISM recusa e sugere a visão esquemática comum;
  \item a síntese para este modo tem \emph{timeout} de 45 segundos.
\end{itemize}

\begin{nota}
Use o PRISM cedo e com frequência: compilar o \code{media_movel} e abrir o diagrama depois de cada mudança de código torna palpável a relação entre construções \cmm{} e custo de hardware --- adicione uma divisão ao programa e veja o divisor aparecer no desenho.
\end{nota}
